\section{Conclusions}
\label{sec:conclusions}

An analytical design-space sweep maps conductance, resistance shares, and six
feasibility criteria for aluminum and PA tubes.
The sweep is anchored to the measured aluminum benchmark from the companion
study~\cite{Henzler2026}; PA points are extrapolated from that basis.
All maps hold at the fixed operating point
($u_{\mathrm{air},\infty}=5\,\mathrm{m\,s^{-1}}$,
$v_{\mathrm{i}}=0.5\,\mathrm{m\,s^{-1}}$) with the VDI G1/G7 correlations; they
are local correlation results for this operating window.
The conclusions are:

\begin{enumerate}

\item In the small-diameter PA minimum-wall interval
  ($d_{\mathrm{o}}=0.3$--$2.0\,\mathrm{mm}$), Fig.~\ref{fig:shares} shows that
  the wall remains a minor resistance term and outer convection remains the
  controlling resistance term.
  The conductivity penalty is governed by the fractional wall-resistance share,
  so the 880-fold conductivity difference to aluminum enters through a small
  resistance fraction.

\item With the six criteria applied under the normal-route tolerance
  ($\Delta t=0.020\,\mathrm{mm}$), the design-boundary map places the
  case-study PA target near $d_{\mathrm{o}}\approx 0.35\,\mathrm{mm}$, where
  the exploratory tube-supply-cost cutoff is more restrictive than the updated
  PA/process capillary limit under the $G=10$,
  $h_{\mathrm{cap}}\le2\,\mathrm{mm}$ case (Fig.~\ref{fig:design_boundary}).
  At sub-millimeter diameters, this interval requires wall thicknesses
  moderately above the process minimum, about
  $0.030$--$0.040\,\mathrm{mm}$, to pass the $6\,\mathrm{bar}$ burst criterion.
  The aluminum minimum-wall option starts at
  $d_{\mathrm{o}}\approx 0.63$--$0.65\,\mathrm{mm}$, because the
  tube-supply-cost cutoff excludes smaller aluminum diameters under the
  benchmark-pitch assumptions.
  Admission of PA without an equal-diameter aluminum comparator remains
  conditional.

\item Because the identified PA interval depends on controlled potting at
  $G=10$, a derated PA burst-strength input, and wall thicknesses above the
  process minimum under the normal $\pm 20\,\mu\mathrm{m}$ extrusion tolerance,
  the practical design limit shifts from thermal resistance to manufacturing,
  pressure integrity, and durability.
  Early prototype builds shifted the failure concern from isolated tube burst
  toward frame stiffness, bundle deformation, and tube--potting integrity.
  The next qualification step is bundle-level structural testing of
  wall-tolerance capability, potting adhesion, pressure integrity, and
  long-term coolant-interface durability.

\end{enumerate}
